\documentclass[11pt]{article}
\usepackage{amsmath,amssymb,amsthm}
\usepackage[margin=1in]{geometry}
\usepackage{enumitem}

\newcommand{\XiBC}{\Xi_{\mathrm{BC}}}
\newcommand{\XiH}{\Xi_{\mathrm{BC,Hecke}}}
\newcommand{\Q}{\mathbb Q}
\newcommand{\Hzeta}{H_{\zeta}^{-}}
\newcommand{\HH}{H_{Q,1}^{-}}
\newcommand{\Yzeta}{Y_{\zeta}^{-}}
\newcommand{\YH}{Y_{Q,1}^{-}}

\title{Step 168: H6 Zeta-Fiber Descent Audit}
\author{RH Track Adequacy Ledger}
\date{}

\begin{document}
\maketitle

\section{Purpose}

This note audits H6, the bridge from the Hecke residual \(\XiH\) at
\(K=\Q\), \(\chi=1\), to the inherited Burnol/Sonine residual \(\XiBC\).
The output is one typed verdict among equality, defect-paid transfer,
non-comparability, or audit indeterminacy.

\section{T1. Typed Bridge Context}

The Hecke fiber under audit is:
\[
  \bigl(\HH,L_{Q,1},D_{Q,1},C_{Q,1}\bigr),
\]
where \(\HH\) is the \(K=\Q\), trivial-character Hilbert carrier,
\(L_{Q,1}\) is its native probe family, \(D_{Q,1}\) is its dissolving probe
family, and \(C_{Q,1}\) is its audit energy.

The Burnol/Sonine side is the inherited pulled-evaluator carrier:
\[
  \bigl(H,L,D,C\bigr),
\]
with parent residual
\[
  \XiBC = B^* \Pi_Y B.
\]

Step 92 records candidate bridge notation
\[
  B:\Yzeta\to \YH,\qquad R:\YH\to \Yzeta,
\]
and the desired bridge pattern
\[
  RB=I,\qquad K_{\zeta}\preceq R K_H R^*+E_{\mathrm{br}}.
\]
For step 168 these are candidate or required bridge slots. They are not
accepted comparison data.

The comparison target is therefore either an exact residual identity
\[
  \XiH|_{K=\Q,\chi=1}=\XiBC,
\]
or a defect inequality transferring the Hecke-fiber conclusion to the
Burnol/Sonine carrier through a positive-budgeted \(E_{\mathrm{br}}\).

\section{T2. Source Audit}

Step 92 supplies a Hecke/idele carrier sketch. It defines a completed
Hecke response space, character readouts, a source-frame program, and
bridge slots \(B\) and \(R\). It classifies the Hecke-to-zeta bridge as a
major required object. The machine records mark the bridge as
``major\_bridge\_required'', ``missing\_key'', or
``enlarged\_carrier\_program'' rather than as accepted comparison data.

Step 167 imports this status. It records:
\begin{itemize}[leftmargin=2em]
\item BR1: \(L_{\Q}(s,1)=\zeta(s)\), accepted only at scalar L-function level.
\item BR2: comparison of Hecke fiber carrier with Burnol/Sonine carrier,
status open external bridge.
\item BR3: step 92 bridge maps and defect inequality, status open external bridge.
\end{itemize}

The adjacent records through steps 30--95 supply framework discipline,
source-frame schemas, and finite-window warnings. They do not supply a
carrier isomorphism \(\HH\simeq H\), probe congruences
\(L_{Q,1}\leftrightarrow L\), dissolving-probe congruences
\(D_{Q,1}\leftrightarrow D\), an audit-energy identity or inequality
between \(C_{Q,1}\) and \(C\), or a zero-ledger congruence beyond scalar
L-function equality.

\section{Verdict}

\[
\boxed{\texttt{bridge\_verdict = non\_comparability}.}
\]

This verdict is ledger-relative: under the inherited records, the Hecke
fiber and Burnol/Sonine carrier are structurally non-comparable for H6.
The verdict does not forbid a future external bridge theorem. It says no
such theorem is present in the inherited record.

\section{T3. Framework-Derived No-Go}

\textbf{No-go statement.}
The Hecke fiber carrier at \(K=\Q\), \(\chi=1\), is structurally
non-comparable to the Burnol/Sonine pulled-evaluator carrier under the
inherited records; consequently, closure of \(\XiH\) does not supply closure
of \(\XiBC\) on the Burnol/Sonine carrier.

\textbf{Scope.}
The no-go applies only to H6 under the step 168 ledger. It does not alter
the Burnol/Sonine cascade and it does not alter H1--H5 on the Hecke side.
It also does not deny the scalar identity \(L_{\Q}(s,1)=\zeta(s)\). Its
point is that scalar identity is not carrier identity.

\textbf{Proof sketch.}
Equality would require accepted \(B,R\) with \(RB=I\) and congruence of
Hilbert carrier, native probes, dissolving probes, audit energy, and zero
ledger. Step 92 does not supply those comparison data. Defect-paid transfer
would require an explicit positive-budgeted \(E_{\mathrm{br}}\) satisfying
the defective bridge theorem hypotheses. Step 92 names a bridge-defect slot
but does not supply positivity, fixed-ledger control, or an adequacy budget.
The only accepted cross-record identity is \(L_{\Q}(s,1)=\zeta(s)\), and
step 167 explicitly retains the no-go that scalar L-function identity is not
carrier identity. Therefore H6 is non-comparable under the inherited records.

\textbf{Nonclaim record.}
This step does not assert exact equality, does not assert a defect-paid
transfer, does not close either parent residual, and does not promote BR2 or
BR3 to accepted bridge sources.

\section{T5. Residual Tree Update}

The Burnol/Sonine tree remains a sibling residual tree with unchanged open
branches:
\begin{itemize}[leftmargin=2em]
\item Branch A: Calkin bridge G2--G5, split external theorem.
\item Branch B: per-zero finite-carrier SL164.1, finite-carrier diagnostic
indeterminate.
\item Branch C: shifted co-Poisson H1--H5, split external theorem.
\end{itemize}

The Hecke tree remains active with H1--H5 unchanged. H6 is updated:
\[
  \text{H6 zeta-fiber descent to }\XiBC:
  \quad \texttt{open\_external\_bridge}
  \longmapsto
  \texttt{non\_comparability}.
\]

\end{document}

